\textbf{Phần 1. Trắc nghiệm lựa chọn}
\setcounterpageref{ex}{0}

\begin{ex} %[2H9N46]
Vật nào sau đây có dạng khối cầu?
\choice
{Hộp sữa bột dạng lon}
{Viên gạch}
{Quyển sách}
{\True Quả bóng đá}
\loigiai{
Quả bóng đá có dạng khối cầu.\\
\textbf{Đáp án của câu hỏi là D}
}
\end{ex}

\begin{ex} %[2H9N46]
Vật nào sau đây có dạng khối trụ?
\choice
{Quả địa cầu}
{\True Lon nước ngọt}
{Hộp phấn viết bảng}
{Quả bóng bàn}
\loigiai{
Lon nước ngọt có dạng khối trụ.\\
\textbf{Đáp án của câu hỏi là B}
}
\end{ex}

\begin{ex} %[2H9H46]
Khối cầu có đặc điểm nào sau đây?
\choice
{Có các mặt phẳng vuông vức}
{Chỉ lăn được về hai phía}
{Có 6 mặt bằng nhau}
{\True Lăn được về mọi phía}
\loigiai{
Khối cầu lăn được về mọi phía.\\
\textbf{Đáp án của câu hỏi là D}
}
\end{ex}

\begin{ex} %[2H9H46]
Khối trụ có 2 đáy là hai hình gì?
\choice
{\True Hình tròn bằng nhau}
{Hình vuông}
{Hình tam giác}
{Hình chữ nhật}
\loigiai{
Khối trụ có 2 đáy là 2 hình tròn bằng nhau.\\
\textbf{Đáp án của câu hỏi là A}
}
\end{ex}

\begin{ex} %[2H9H46]
Khi đặt khối trụ đứng trên mặt bàn, khối trụ có thể
\choice
{Lăn tròn}
{Tự quay tròn}
{\True Trượt hoặc đặt chồng lên nhau}
{Lăn về mọi phía}
\loigiai{
Khi đặt đứng thì khối trụ trượt được và xếp chồng lên nhau được.\\
\textbf{Đáp án của câu hỏi là C}
}
\end{ex}

\begin{ex} %[2H9H46]
Khi đặt khối trụ nằm ngang trên mặt bàn, khối trụ có thể
\choice
{Đặt chồng vững chắc}
{\True Lăn được}
{Không di chuyển được}
{Lăn về mọi phía như khối cầu}
\loigiai{
Khi đặt nằm ngang khối trụ có thể lăn được.\\
\textbf{Đáp án của câu hỏi là B}
}
\end{ex}

\begin{ex} %[2H9V46]
Trong các vật sau: quả cam, hộp chè hình trụ, viên bi, hộp bánh hình lập phương. Có bao nhiêu vật có dạng khối cầu?
\choice
{$1$ vật}
{$3$ vật}
{\True $2$ vật (quả cam, viên bi)}
{$4$ vật}
\loigiai{
Có 2 vật: quả cam và viên bi.\\
\textbf{Đáp án của câu hỏi là C}
}
\end{ex}

\begin{ex} %[2H9V46]
Trong các vật sau: lon sữa đặc, ống bơ đong gạo, khối ru-bích, khúc gỗ tròn. Có bao nhiêu vật có dạng khối trụ?
\choice
{\True $3$ vật (lon sữa, ống bơ, khúc gỗ tròn)}
{$2$ vật}
{$4$ vật}
{$1$ vật}
\loigiai{
Có 3 vật có dạng khối trụ.\\
\textbf{Đáp án của câu hỏi là A}
}
\end{ex}

\begin{ex} %[2H9V46]
Để xếp chồng được lên nhau một cách vững chắc nhất, người ta nên dùng khối nào?
\choice
{Khối cầu}
{Quả bóng}
{Viên bi}
{\True Khối trụ (đặt đứng) hoặc khối lập phương}
\loigiai{
Khối trụ đặt đứng hoặc lập phương có mặt phẳng để xếp chồng.\\
\textbf{Đáp án của câu hỏi là D}
}
\end{ex}

\begin{ex} %[2H9V46]
Bạn An có 4 khối cầu và 5 khối trụ. An có tất cả số khối hình là
\choice
{$8$ khối}
{$10$ khối}
{\True $9$ khối}
{$11$ khối}
\loigiai{
$4 + 5 = 9$ khối.\\
\textbf{Đáp án của câu hỏi là C}
}
\end{ex}

\begin{ex} %[2H9V46]
Vật nào sau đây KHÔNG có dạng khối cầu?
\choice
{Quả bóng rổ}
{\True Chiếc nón lá}
{Hòn bi ve}
{Quả bóng bàn}
\loigiai{
Chiếc nón lá có dạng hình nón (không phải khối cầu).\\
\textbf{Đáp án của câu hỏi là B}
}
\end{ex}

\begin{ex} %[2H9C46]
Xếp 3 khối trụ giống nhau chồng lên nhau theo phương thẳng đứng. Hình thu được vẫn là một
\choice
{\True Khối trụ cao hơn}
{Khối cầu}
{Khối lập phương}
{Khối chữ nhật}
\loigiai{
Xếp chồng các khối trụ cùng đáy tạo thành 1 khối trụ cao hơn.\\
\textbf{Đáp án của câu hỏi là A}
}
\end{ex}

\textbf{Phần 2. Trắc nghiệm đúng sai}
\setcounterpageref{ex}{0}

\begin{ex} %[2H9H46]
Xét tính đúng sai về hình khối:
\choiceTF
{\True Khối cầu không có cạnh và không có đỉnh.}
{\True Khối trụ có hai mặt đáy phẳng hình tròn.}
{Khối cầu có thể xếp chồng lên nhau một cách dễ dàng.}
{\True Khối cầu dễ bị lăn nên rất khó xếp chồng trực tiếp lên nhau.}
\loigiai{
\begin{itemize}
    \item Khối cầu không có cạnh, đỉnh: Đúng.
    \item Khối trụ có 2 đáy tròn: Đúng.
    \item Khối cầu trơn lăn khó xếp chồng. Ý 3 Sai, Ý 4 Đúng.
\end{itemize}
\begin{itemize}
    \item a) \textbf{Đúng}.
    \item b) \textbf{Đúng}.
    \item c) \textbf{Sai}.
    \item d) \textbf{Đúng}.
\end{itemize}
}
\end{ex}

\begin{ex} %[2H9V46]
Cho tập hợp các đồ vật: quả bóng đá, lon cocacola, bình thủy tinh hình trụ, viên bi ve. Xét tính đúng sai:
\choiceTF
{\True Có 2 vật có dạng khối cầu (quả bóng, viên bi).}
{\True Có 2 vật có dạng khối trụ (lon cocacola, bình thủy tinh).}
{Tất cả các vật trên đều có dạng khối cầu.}
{\True Các vật trên gồm 2 khối cầu và 2 khối trụ.}
\loigiai{
\begin{itemize}
    \item 2 khối cầu: Đúng.
    \item 2 khối trụ: Đúng.
    \item Không phải tất cả là khối cầu. Ý 3 Sai, Ý 4 Đúng.
\end{itemize}
\begin{itemize}
    \item a) \textbf{Đúng}.
    \item b) \textbf{Đúng}.
    \item c) \textbf{Sai}.
    \item d) \textbf{Đúng}.
\end{itemize}
}
\end{ex}

\textbf{Phần 3. Trả lời ngắn}
\setcounterpageref{ex}{0}

\begin{ex} %[2H9H46]
Khối trụ có mấy mặt đáy hình tròn?
\shortans{2}
\loigiai{
\textbf{Đáp số:} $2$\\
2 mặt đáy.
}
\end{ex}

\begin{ex} %[2H9H46]
Quả bóng đá có dạng khối gì? (Viết: khối cầu)
\shortans{khối cầu}
\loigiai{
\textbf{Đáp số:} khối cầu\\
Khối cầu.
}
\end{ex}

\begin{ex} %[2H9V46]
Lon nước ngọt dạng hình gì? (Viết: khối trụ)
\shortans{khối trụ}
\loigiai{
\textbf{Đáp số:} khối trụ\\
Khối trụ.
}
\end{ex}

\begin{ex} %[2H9V46]
Có 6 khối cầu và 8 khối trụ. Hỏi có tất cả bao nhiêu khối hình?
\shortans{14}
\loigiai{
\textbf{Đáp số:} $14$\\
$6 + 8 = 14$ khối.
}
\end{ex}

\begin{ex} %[2H9V46]
Có 15 khối trụ, chia đều thành 5 cột bằng nhau. Hỏi mỗi cột có bao nhiêu khối trụ?
\shortans{3}
\loigiai{
\textbf{Đáp số:} $3$\\
$15 : 5 = 3$ khối.
}
\end{ex}

\begin{ex} %[2H9C46]
Bé Nam dùng 2 khối lập phương và 3 khối trụ để xếp một tòa tháp. Hỏi tòa tháp được tạo từ bao nhiêu khối hình?
\shortans{5}
\loigiai{
\textbf{Đáp số:} $5$\\
$2 + 3 = 5$ khối.
}
\end{ex}
